% sedfit User Manual
% Build with:  pdflatex sedfit_user_manual.tex   (run twice for the TOC)
%
% The format block below is shared verbatim with the sedphot manual.

\documentclass[11pt]{article}

\newcommand{\manualname}{sedfit}
\newcommand{\manualtagline}{Photometry in, redshifts out}
\newcommand{\manualversion}{0.1.0}

% ----------------------------------------------------------------
% FORMAT. This block is shared verbatim with the sibling user
% manual in the other repository so the two are typeset
% identically. Change it in both or in neither.
% ----------------------------------------------------------------
\usepackage[T1]{fontenc}
\usepackage{lmodern}
\usepackage{textcomp}
\usepackage{csquotes}

% Straight quotes in monospace, curly apostrophes in prose, so a command
% copied out of a code block pastes back into a shell unchanged.
% Typewriter faces have zero interword stretch; that is the test.
%
% Section titles must use \textquoteright rather than a bare ' --- the
% kernel's math-active ' wins when a title is written to the .toc, and
% the entry comes out as a math prime.
\begingroup
  \catcode`\'=\active
  \protected\gdef'{\ifdim\fontdimen3\font=0pt
                      \textquotesingle\else\char39\relax\fi}
\endgroup
\AtBeginDocument{\catcode`\'=\active}
\usepackage[a4paper,margin=1in]{geometry}
\usepackage{fancyhdr}
\usepackage[dvipsnames]{xcolor}
\usepackage{booktabs}
\usepackage{longtable}
\usepackage{array}
\usepackage{enumitem}
\usepackage{titlesec}
\usepackage[most]{tcolorbox}
\usepackage[hidelinks]{hyperref}
\usepackage{url}

% --- Palette ---
\definecolor{slate}{RGB}{47,72,88}          % headings
\definecolor{boxrule}{RGB}{201,204,209}     % code-block border
\definecolor{boxfill}{RGB}{250,250,251}     % code-block fill
\definecolor{noteaccent}{RGB}{180,95,20}
\definecolor{warnaccent}{RGB}{160,30,30}

% --- Headings ---
\titleformat{\section}{\Large\bfseries\color{slate}}{\thesection}{1em}{}
\titleformat{\subsection}{\large\bfseries\color{slate}}{\thesubsection}{1em}{}

% --- Running head ---
\pagestyle{fancy}
\fancyhf{}
\lhead{\small\manualname\ user manual}
\rhead{\small\thepage}
\renewcommand{\headrulewidth}{0.4pt}

% --- Blocks ---
% A shell/code block: light fill, thin grey border, monospaced.
\newtcolorbox{cmd}{
  colback=boxfill, colframe=boxrule, boxrule=0.4pt,
  left=6pt, right=6pt, top=5pt, bottom=5pt,
  before skip=8pt, after skip=8pt, breakable
}

\newtcolorbox{notebox}[1][Note]{
  breakable, colback=noteaccent!5, colframe=noteaccent,
  boxrule=0.8pt, title=#1, fonttitle=\bfseries\small
}

\newtcolorbox{warnbox}[1][Careful]{
  breakable, colback=warnaccent!5, colframe=warnaccent,
  boxrule=0.8pt, title=#1, fonttitle=\bfseries\small
}

% One-line command.
\newcommand{\cmdline}[1]{\begin{cmd}\ttfamily #1\end{cmd}}

\setlist[itemize]{itemsep=2pt, topsep=4pt}
\setlist[enumerate]{itemsep=2pt, topsep=4pt}

% Table helpers: booktabs rules, wrapped description columns.
\newcolumntype{L}[1]{>{\raggedright\arraybackslash}p{#1}}
\newcolumntype{T}[1]{>{\ttfamily\raggedright\arraybackslash}p{#1}}

% --- Title page ---
\newcommand{\manualtitlepage}{%
  \thispagestyle{empty}
  \vspace*{4cm}
  \begin{center}
    {\Huge\bfseries\manualname}\\[0.7cm]
    {\Large User Manual}\\[1.3cm]
    {\normalsize\manualtagline}
  \end{center}
  \vfill
  \begin{center}\normalsize Version \manualversion\end{center}
  \newpage
  \tableofcontents
  \newpage
}

\begin{document}

\manualtitlepage

% ==========================================================
\section{What this does}
% ==========================================================

\texttt{sedfit} takes photometry you already have for a set of galaxies
and fits each one's spectral energy distribution to estimate its
redshift.

You give it three things:

\begin{enumerate}
  \item a \textbf{sample catalog} --- a spreadsheet listing your
        galaxies, one row each;
  \item a \textbf{campaign config} --- a settings file saying which
        photometry files to expect and how to combine them;
  \item a \textbf{fit config} --- a settings file saying how to fit.
\end{enumerate}

It then does three things:

\begin{enumerate}
  \item \textbf{Builds} one combined photometry table per galaxy, putting
        every instrument on a common flux scale.
  \item \textbf{Fits} that table with one of two engines and estimates a
        redshift with uncertainties.
  \item \textbf{Records} everything --- every input file is hashed, every
        setting is saved next to the results, and every run is logged in
        one central file.
\end{enumerate}

Fluxes are in microjansky on the AB system throughout. Everything for one
galaxy lands under that galaxy's own directory.

% ==========================================================
\section{Installing}
% ==========================================================

These instructions assume the machine has nothing set up yet. Commands
are the same on all three operating systems unless stated otherwise.

\subsection{Get Python}

\texttt{sedfit} needs \textbf{Python 3.11 or newer} plus a scientific
stack (NumPy, SciPy, Astropy, Pandas, Matplotlib).

Use \textbf{Miniconda}. Conda installs Python and the scientific
libraries as pre-built binaries, which avoids the compiler toolchain some
of these packages otherwise need, and it keeps each project's libraries
separate so one project cannot break another.

\textbf{Windows}
\begin{enumerate}
  \item Go to \url{https://www.anaconda.com/download/success} and
        download the \textbf{Miniconda Windows 64-bit installer}
        (\texttt{.exe}).
  \item Run the installer and accept the defaults.
  \item From the Start menu, open \textbf{Anaconda Prompt}. Every command
        in this manual is typed there.
\end{enumerate}

\textbf{macOS}
\begin{enumerate}
  \item Download the \textbf{Miniconda macOS installer} for your chip
        (Apple silicon = \texttt{arm64}, Intel = \texttt{x86\_64}).
  \item Run it, then open \textbf{Terminal}.
\end{enumerate}

\textbf{Linux}
\begin{enumerate}
  \item Download the Miniconda Linux \texttt{x86\_64} installer script.
  \item Run \texttt{bash Miniconda3-latest-Linux-x86\_64.sh}, then open a
        new terminal.
\end{enumerate}

Check that conda is available:

\cmdline{conda -{}-version}

You should see something like \texttt{conda 24.9.2}.

\subsection{Get the code}

The source lives at \url{https://github.com/chrishopp12/sed-fitting}.

\textbf{With git.} On Windows, install Git for Windows from
\url{https://git-scm.com/download/win} first; macOS and Linux usually
have git already. Then:

\cmdline{git clone https://github.com/chrishopp12/sed-fitting.git}

This creates a folder named \texttt{sed-fitting}. Move into it:

\cmdline{cd sed-fitting}

\textbf{Without git.} Open the page in a browser, click the green
\textbf{Code} button, choose \textbf{Download ZIP}, and unzip it. You get
a folder named \texttt{sed-fitting-main}. Move into it:

\cmdline{cd sed-fitting-main}

You are in the right folder if it contains a file named
\texttt{pyproject.toml}.

\begin{notebox}[The download is about 20 MB and unpacks to about 70 MB]
Most of that is the packaged template libraries --- the galaxy spectra
the fit compares your photometry against. They ship with the code so a
fit is reproducible without downloading anything else.
\end{notebox}

\subsection{Create an environment and install}

\textbf{Conda route (recommended, all platforms).}

\cmdline{conda create -n sedfit python=3.11}

Answer \texttt{y} when it asks. Then activate it:

\cmdline{conda activate sedfit}

Your prompt now starts with \texttt{(sedfit)}. Install the package:

\cmdline{pip install -e .}

That reads \texttt{pyproject.toml} and pulls in every dependency.

If a dependency fails to build, install the scientific stack from
conda-forge first and then install \texttt{sedfit} without re-resolving:

\cmdline{conda install -c conda-forge numpy scipy pandas astropy matplotlib}

\cmdline{pip install -e . -{}-no-deps}

\textbf{pip / venv route (no conda).} Only use this if Python 3.11+ is
already installed and on your PATH.

Windows:

\cmdline{python -m venv .venv}

\cmdline{.venv\textbackslash Scripts\textbackslash activate}

macOS / Linux:

\cmdline{python3 -m venv .venv}

\cmdline{source .venv/bin/activate}

Then, on any platform:

\cmdline{pip install -e .}

\subsection{What \enquote{editable} means}

\texttt{pip install -e .} installs the package \textbf{in place}: Python
runs the source files in the folder you downloaded, rather than a copy
hidden in the environment. If you \texttt{git pull} a newer version, the
change takes effect immediately with no reinstall.

\subsection{Optional extras}

The base install runs the assembly path and the quick engine --- enough
for your first run. Four extras exist:

\cmdline{pip install -e ".[eazy]"\qquad\# the official eazy-py engine}

\cmdline{pip install -e ".[prospector]"\qquad\# Prospector + FSPS}

\cmdline{pip install -e ".[test]"\qquad\# pytest, for the check below}

\cmdline{pip install -e ".[dered]"\qquad\# astroquery, for -{}-deredden}

The last two are easy to overlook: \texttt{pytest} is not in the base
install, so \texttt{python -m pytest} below fails without
\texttt{".[test]"}, and the automatic E(B--V) lookup that
\texttt{-{}-deredden} performs needs \texttt{astroquery}. You can install
several at once: \texttt{pip install -e ".[eazy,test,dered]"}.

\begin{warnbox}[Prospector on Windows]
Prospector needs FSPS, which compiles from Fortran source and is
difficult to build on Windows. The base install and both eazy engines run
on any platform; if you need Prospector, use Linux or macOS.
\end{warnbox}

Prospector also needs a copy of the FSPS data files and an environment
variable pointing at them, set before every run:

\cmdline{git clone https://github.com/cconroy20/fsps \textasciitilde/fsps}

\cmdline{export SPS\_HOME=\textasciitilde/fsps}

\subsection{Choosing a stellar library: MILES or C3K}

Prospector builds a galaxy's light out of model stellar spectra, and
FSPS ships two libraries to draw them from: \textbf{MILES} (measured
from real stars) and \textbf{C3K} (computed from stellar atmosphere
models). They do not give the same answer, so every Prospector config
has to say which one it used --- that is the \texttt{stellar\_library}
field, either \texttt{"miles"} or \texttt{"c3k"}.

The awkward part is that the library is fixed when FSPS is
\textbf{compiled}, not when it runs. No setting switches between them.
You need a separate environment for each, and plain
\texttt{pip install fsps} gives you the MILES build.

\textbf{MILES} works out of the box:

\begin{cmd}\ttfamily
conda env create -f envs/sedfit\_miles.yml\\
conda activate sedfit\_miles\\
pip install -e .
\end{cmd}

\begin{warnbox}[The .yml files do not install sedfit itself]
They build the environment around it --- Python, the scientific stack,
Prospector, FSPS --- but not the package. Creating an environment and
stopping there gets you \texttt{sedfit: command not found}. Always
follow \texttt{conda env create} with \texttt{conda activate} and
\texttt{pip install -e .} from inside the repository folder.
\end{warnbox}

\textbf{C3K} takes one extra step. Creating the environment is not
enough: its \texttt{.yml} installs the same pre-built package everyone
else gets, which is MILES. You have to rebuild FSPS afterwards, and that
compiles from Fortran source, so the environment needs a compiler first.

\begin{cmd}\ttfamily
conda env create -f envs/sedfit\_c3k.yml\\
conda activate sedfit\_c3k\\
conda install -c conda-forge compilers\\
FFLAGS="-DMILES=0 -DC3K=1" pip install -{}-force-reinstall \textbackslash\\
\hspace*{4em}-{}-no-cache-dir -{}-no-binary fsps fsps==0.4.7\\
pip install -e .
\end{cmd}

That takes a few minutes. macOS ships no Fortran compiler, which is why
the \texttt{compilers} line comes first.

\textbf{Check what a given environment actually has:}

\begin{cmd}\ttfamily
python -c "import fsps; \textbackslash\\
\hspace*{2em}print(fsps.StellarPopulation(zcontinuous=1).libraries)"
\end{cmd}

The middle entry is the spectral library:

\begin{cmd}\ttfamily
sedfit\_miles~~->~~(b'mist', b'miles', b'DL07')\\
sedfit\_c3k~~~~->~~(b'mist', b'c3k\_a', b'DL07')
\end{cmd}

If the C3K environment still reports \texttt{miles}, the rebuild did not
take. Run it again and watch for compiler errors --- a failed build
leaves the previous install in place rather than removing it. If it will
not go through at all and you have another environment that already
reports \texttt{c3k\_a}, clone that instead and install into the clone:

\begin{cmd}\ttfamily
conda create -n sedfit\_c3k -{}-clone <the working env>\\
conda activate sedfit\_c3k\\
pip install -e .
\end{cmd}

\begin{notebox}[The fit checks this for you]
Every Prospector fit compares the config's \texttt{stellar\_library}
against the library the live FSPS build reports, before it does any
work, and stops if they disagree:

\smallskip
\begin{cmd}\ttfamily\small
stellar\_library mismatch: the config declares 'c3k' (FSPS stamp\\
'c3k\_a') but the live build reports 'miles'; run in the matching\\
sedfit env
\end{cmd}
\smallskip

Activating the wrong environment therefore costs you a clear error, not
a quietly wrong answer. Nothing is written when this fires, so there is
no half-finished run to clean up.
\end{notebox}

Results from the two libraries are not interchangeable. Compare fits
within one library, or treat the difference between them as one of the
things you are measuring.

\subsection{Verify the install}

\cmdline{sedfit -{}-help}

You should see a list of seven verbs: \texttt{roster}, \texttt{build},
\texttt{fit}, \texttt{batch}, \texttt{run}, \texttt{plot},
\texttt{manifest}. Then run the test suite (this needs the
\texttt{test} extra above, which installs \texttt{pytest}):

\cmdline{python -m pytest tests/}

Everything should pass. Tests for an engine you did not install skip
themselves rather than failing, so a base install reports a couple of
skips and no failures.

% ==========================================================
\section{Before you start}
% ==========================================================

\texttt{sedfit} does not download photometry. It expects the photometry
files to be on disk already, laid out one directory per galaxy.

\subsection{The directory layout it expects}

\begin{cmd}\ttfamily
data\_root/\\
\hspace*{2em}gal\_0001/\\
\hspace*{4em}Photometry/\\
\hspace*{6em}gal\_0001\_catalog.csv\\
\hspace*{6em}gal\_0001\_measured.csv\\
\hspace*{6em}SPHEREx/\\
\hspace*{8em}table\_photometry.csv\\
\hspace*{2em}gal\_0002/\\
\hspace*{4em}Photometry/\\
\hspace*{6em}...
\end{cmd}

You choose where \texttt{data\_root} is and what the galaxy directories
are called; the campaign config points at them. Names like
\texttt{Photometry/} are conventions the campaign config declares, not
rules baked into the code.

\subsection{What a photometry file must contain}

Every source CSV needs four columns:

\begin{center}
\begin{tabular}{ll}
\toprule
\textbf{Column} & \textbf{Meaning} \\
\midrule
\texttt{band} & the filter name, e.g. \texttt{Legacy\_g} \\
\texttt{flux\_uJy} & flux in microjansky \\
\texttt{flux\_err\_uJy} & its uncertainty \\
\texttt{source} & where the measurement came from \\
\bottomrule
\end{tabular}
\end{center}

The \texttt{source} column matters more than it looks. One file can hold
measurements from several archives, and \texttt{sedfit} tells them apart
by the start of this string --- \texttt{Legacy\_DR9},
\texttt{SDSS\_DR17\_cModel}, and so on. It is how the same band can come
from two different places in one file.

Three more columns are used if you supply them, and ignored if you do
not: \texttt{target\_ra} and \texttt{target\_dec}, which are checked
against the galaxy's catalog position and must agree within
0.5\,arcsec; and \texttt{flags}, which becomes the fit's quality-gate
column for aperture measurements. Any other column is left alone.

\begin{notebox}[Where the legal band names live]
\texttt{sedfit/data/registry.json} lists every band the package knows
--- 37 of them across CFHT, GALEX, J-PLUS, Legacy, Pan-STARRS, SDSS,
SPHEREx and WISE. Open it to check a spelling; an unrecognized name is
a hard error naming what it did not recognize. Watch for the ones that
do not read the way you would guess: Legacy has \texttt{g r i z} and no
\texttt{u}, and the WISE bands are \texttt{WISE\_W1}--\texttt{W4} while
the \emph{provider} tokens for them are \texttt{unwise} and
\texttt{allwise}. SPHEREx channels are the exception --- they are built
per galaxy and are not listed in the registry.
\end{notebox}

\subsection{The three files you write}

\textbf{1. The sample catalog} (a CSV, one row per galaxy). This is the
file you edit as the sample grows.

\begin{cmd}\ttfamily
name,ra\_deg,dec\_deg,dir,z\_ref\_kind,z\_ref\\
gal\_0001,217.51,56.99,gal\_0001,spec,0.1043\\
gal\_0002,217.62,57.02,gal\_0002,spec,0.1058\\
gal\_0003,217.44,56.95,gal\_0003,reference,
\end{cmd}

\begin{center}
\begin{tabular}{ll}
\toprule
\textbf{Column} & \textbf{Meaning} \\
\midrule
\texttt{name} & the galaxy's name; must be unique \\
\texttt{ra\_deg}, \texttt{dec\_deg} & position in degrees \\
\texttt{dir} & its directory under \texttt{data\_root} (optional; defaults to the name) \\
\texttt{z\_ref\_kind} & \texttt{spec}, \texttt{phot}, or \texttt{reference} \\
\texttt{z\_ref} & the known redshift; \textbf{leave blank} for \texttt{reference} \\
\texttt{label} & the filename stem (optional; defaults from the name) \\
\bottomrule
\end{tabular}
\end{center}

\texttt{z\_ref\_kind} says where the comparison redshift comes from.
Use \texttt{spec} or \texttt{phot} when the galaxy has its own measured
value, and \texttt{reference} when it should adopt the sample-wide
redshift set in the campaign config. Columns \texttt{sedfit} does not
recognize are ignored and reported, so you can keep your own notes in the
same file.

\textbf{2. The campaign config} (\texttt{campaign.json}). Everything
shared across galaxies. A minimal one:

\begin{cmd}\ttfamily
\{\\
\hspace*{2em}"schema\_version": 1,\\
\hspace*{2em}"sample": "my\_sample",\\
\hspace*{2em}"data\_root": "data",\\
\hspace*{2em}"reference\_redshift": 0.1050,\\
\hspace*{2em}"position\_authority": "Legacy DR9 positions",\\
\hspace*{2em}"sources": \{\\
\hspace*{4em}"legacy": \{\\
\hspace*{6em}"path": "Photometry/\{label\}\_catalog.csv",\\
\hspace*{6em}"bands": ["Legacy\_g", "Legacy\_r", "Legacy\_z"],\\
\hspace*{6em}"kind": "catalog",\\
\hspace*{6em}"provider": "legacy"\\
\hspace*{4em}\}\\
\hspace*{2em}\},\\
\hspace*{2em}"recipes": \{\\
\hspace*{4em}"plain": \{\\
\hspace*{6em}"reference": "none",\\
\hspace*{6em}"sources": [\{"source": "legacy", "role": "float"\}],\\
\hspace*{6em}"spherex": null\\
\hspace*{4em}\}\\
\hspace*{2em}\}\\
\}
\end{cmd}

\texttt{data\_root} is relative to the campaign file --- so this campaign
sits next to a \texttt{data/} folder holding the galaxy directories. A
\emph{source} is one photometry file and the bands you expect from it;
\texttt{\{label\}} is replaced with each galaxy's filename stem.
\texttt{provider} must be one of \texttt{legacy}, \texttt{unwise},
\texttt{allwise}, \texttt{galex}, \texttt{sdss}, \texttt{panstarrs},
\texttt{jplus}, \texttt{hst}, \texttt{aperture}, \texttt{sersic} --- each
maps to the prefix its rows must carry in the \texttt{source} column
(\texttt{legacy} $\rightarrow$ \texttt{Legacy\_}, \texttt{unwise}
$\rightarrow$ \texttt{unWISE\_Legacy\_}, and so on). A \emph{recipe} is one way
to combine those sources into a table --- you can define several and run
them all. \texttt{"spherex": null} in a recipe means \emph{deliberately
leave the spectrum out of this one}, which is what you want when the
recipe is broadband-only.

\begin{warnbox}[reference\_redshift is required as soon as you use it]
The moment one catalog row says \texttt{z\_ref\_kind: reference}, the
campaign must declare \texttt{reference\_redshift} --- that is where the
redshift comes from. \texttt{sedfit roster} does \textbf{not} check
this: it writes a clean roster, and the \emph{next} command fails with
\texttt{z\_ref\_kind 'reference' requires the roster to declare
reference\_redshift}. If a verb fails that way right after a successful
roster, this is why.
\end{warnbox}

\textbf{3. The fit config} (e.g.\ \texttt{quick.json}). How to fit:

\begin{cmd}\ttfamily
\{\\
\hspace*{2em}"schema\_version": 2,\\
\hspace*{2em}"backend": "eazy",\\
\hspace*{2em}"name": "quick\_v1",\\
\hspace*{2em}"err\_floor": 0.05,\\
\hspace*{2em}"min\_snr\_broadband": 2.0,\\
\hspace*{2em}"min\_valid\_bands": 5,\\
\hspace*{2em}"eazy": \{\\
\hspace*{4em}"engine": "quick",\\
\hspace*{4em}"z\_min": 0.0,\\
\hspace*{4em}"z\_max": 1.0,\\
\hspace*{4em}"z\_step": 0.001,\\
\hspace*{4em}"templates": "brown14\_vac\_cosmos160",\\
\hspace*{4em}"template\_pattern": "*.dat"\\
\hspace*{2em}\}\\
\}
\end{cmd}

\begin{notebox}[You must name a template set]
There is no default. The template set is the single biggest influence on
the redshift you get out, so the code makes you choose it deliberately.
Name any of the sets shipped in \texttt{sedfit/data/templates/} ---
\texttt{brown14\_vac\_cosmos160} is the recommended one --- or give a path
to your own directory of spectra.
\end{notebox}

\begin{warnbox}[And set template\_pattern with it]
\texttt{template\_pattern} says which files in that directory are
spectra, and it defaults to \texttt{*\_spec.dat}. The recommended set
mixes two naming conventions: 129 Brown+14 spectra named
\texttt{*\_spec.dat} and 31 COSMOS templates named differently. So the
default silently fits \textbf{129 of the 160} templates. Set
\texttt{"template\_pattern": "*.dat"} to use the whole basis. The code
warns when a pattern selects only part of a set, but that warning is
easy to lose in a batch log.
\end{warnbox}

\begin{warnbox}[Always set the redshift range]
\texttt{z\_min} and \texttt{z\_max} default to \textbf{0.05} and
\textbf{0.16} --- a narrow window left over from the survey this package
was written for. A config that omits them fits that window and says
nothing about it, so every galaxy comes back between 0.05 and 0.16
whether or not it belongs there. Set them for your own sample, as the
example above does.
\end{warnbox}

% ==========================================================
\section{The verbs}
% ==========================================================

\begin{center}
\begin{tabular}{>{\ttfamily}l p{10.2cm}}
\toprule
\textbf{\normalfont Verb} & \textbf{What it does} \\
\midrule
roster   & Reads your sample catalog and campaign config, checks what is
           actually on disk, and writes \texttt{roster.json} \\
build    & Combines each galaxy's photometry into one fit-ready table \\
fit      & Fits one galaxy with one recipe \\
batch    & Fits every galaxy, in parallel, and writes a report \\
run      & Does \texttt{roster}, \texttt{build} and \texttt{batch} in one command \\
plot     & Redraws the figures for a run you already have \\
manifest & Summarizes every run recorded so far \\
\bottomrule
\end{tabular}
\end{center}

The normal working order is \texttt{roster} $\rightarrow$ \texttt{build}
$\rightarrow$ \texttt{batch}. The \texttt{run} verb chains those three,
and is what you want most of the time.

\texttt{run} always rebuilds every table and always resumes, and neither
is overridable --- it has no \texttt{-{}-build}, \texttt{-{}-no-resume}
or \texttt{-{}-dry-run}. Reach for \texttt{roster} and \texttt{batch}
separately when you need either off.

A command that fails tells the shell so: \texttt{batch} and \texttt{run}
exit nonzero if any job failed, and \texttt{manifest} exits nonzero if
anything it checked is missing. That is what a script should test.

Every verb explains itself:

\cmdline{sedfit -{}-help \qquad sedfit build -{}-help}

% ==========================================================
\section{Your first run}
% ==========================================================

Put your sample catalog and campaign config next to each other, then:

\textbf{Step 1 --- check what the code can see.} This writes nothing:

\begin{cmd}\ttfamily
sedfit roster -{}-catalog sample.csv -{}-campaign campaign.json -{}-dry-run
\end{cmd}

You get a line per galaxy-and-source that did not work out:

\begin{cmd}\ttfamily
sample.csv: 3 rows; using name, ra\_deg, dec\_deg, dir, z\_ref\_kind, z\_ref\\
3 catalog rows -> 3 targets\\
\hspace*{2em}1 dropped, 8 kept\\
\hspace*{2em}dropped~~gal\_0003/legacy: Photometry/gal\_0003\_catalog.csv: absent\\
dry run: nothing written
\end{cmd}

Read that carefully before going on --- it is telling you which files it
found and which it did not.

\textbf{Step 2 --- write the roster.} Drop \texttt{-{}-dry-run}:

\begin{cmd}\ttfamily
sedfit roster -{}-catalog sample.csv -{}-campaign campaign.json
\end{cmd}

This writes \texttt{roster.json} next to the campaign file, plus
\texttt{roster.json.report.csv}, which has one row per galaxy-and-source
with the reason it was kept, partially kept, or dropped.

\begin{warnbox}[Never edit roster.json by hand]
It is generated. If something in it is wrong, fix the sample catalog or
the campaign config and run \texttt{sedfit roster} again --- it takes
seconds. Hand edits are silently overwritten the next time anyone
regenerates it.
\end{warnbox}

\textbf{Step 3 --- build and fit everything:}

\begin{cmd}\ttfamily
sedfit batch -{}-roster roster.json -{}-config quick.json -{}-build
\end{cmd}

You get a line per job as it finishes:

\begin{cmd}\ttfamily
6 jobs (eazy), 3 worker(s); 0 pair(s) not applicable\\
{[1/6]} ok\hspace*{6em}gal\_0001/plain z50=0.1043\\
{[2/6]} ok\hspace*{6em}gal\_0002/plain z50=0.1061\\
...\\
6 ok, 0 failed, 0 skipped -> /path/to/batch\_report.csv
\end{cmd}

\textbf{Or do all three steps at once:}

\begin{cmd}\ttfamily
sedfit run -{}-catalog sample.csv -{}-campaign campaign.json \textbackslash\\
\hspace*{4em}-{}-config quick.json
\end{cmd}

% ==========================================================
\section{Common tasks}
% ==========================================================

\subsection{A note on long commands}

A few commands below are too wide for one line and are shown broken
across two, with a trailing \texttt{\textbackslash}. That is the
continuation character for bash and zsh (macOS, Linux).

\begin{center}
\begin{tabular}{L{6.0cm} L{5.5cm}}
\toprule
\textbf{Shell} & \textbf{Continuation character} \\
\midrule
macOS / Linux (bash, zsh) & \texttt{\textbackslash} \\
Windows Command Prompt / Anaconda Prompt & \texttt{\textasciicircum} \\
Windows PowerShell & \texttt{\textasciigrave}\ (backtick) \\
\bottomrule
\end{tabular}
\end{center}

On Windows, either swap the trailing \texttt{\textbackslash} for the
right character, or --- simplest and always correct --- delete the line
break and paste the whole thing as one long line.

\subsection{Add galaxies to the sample}

Add rows to your sample catalog, then regenerate and re-run:

\begin{cmd}\ttfamily
sedfit roster -{}-catalog sample.csv -{}-campaign campaign.json\\
sedfit batch -{}-roster roster.json -{}-config quick.json -{}-build
\end{cmd}

Galaxies that were already fitted are skipped automatically, so only the
new ones cost time. This is the everyday loop --- unless
\texttt{-{}-build} produces a different table than last time, which
changes the run fingerprint and makes it a new fit.

\subsection{Fit one galaxy}

\begin{cmd}\ttfamily
sedfit fit -{}-roster roster.json -{}-target gal\_0001 \textbackslash\\
\hspace*{4em}-{}-recipe plain -{}-config quick.json
\end{cmd}

\subsection{Check a command before running it}

\texttt{-{}-dry-run} validates everything and writes nothing. It exists
on \texttt{roster}, \texttt{build} and \texttt{fit}; \texttt{batch} and
\texttt{run} have no dry run, so validate one galaxy with \texttt{fit
-{}-dry-run} before launching a sweep. Use it whenever you have changed
a config:

\begin{cmd}\ttfamily
sedfit fit -{}-roster roster.json -{}-target gal\_0001 \textbackslash\\
\hspace*{4em}-{}-recipe plain -{}-config quick.json -{}-dry-run
\end{cmd}

\subsection{Build tables without fitting}

\begin{cmd}\ttfamily
sedfit build -{}-roster roster.json
\end{cmd}

Add \texttt{-{}-target} or \texttt{-{}-recipe} to narrow it down.

\subsection{Run only some galaxies}

Put the names in a CSV with a \texttt{name} column --- your sample
catalog works directly --- and pass it:

\begin{cmd}\ttfamily
sedfit batch -{}-roster roster.json -{}-config quick.json \textbackslash\\
\hspace*{4em}-{}-targets subset.csv
\end{cmd}

Or a single one with \texttt{-{}-target gal\_0001}.

\subsection{Run only one recipe}

\begin{cmd}\ttfamily
sedfit batch -{}-roster roster.json -{}-config quick.json -{}-recipe plain
\end{cmd}

\subsection{Correct for Milky Way dust}

\begin{cmd}\ttfamily
sedfit batch -{}-roster roster.json -{}-config quick.json \textbackslash\\
\hspace*{4em}-{}-build -{}-deredden
\end{cmd}

This looks up the extinction at each galaxy's position and writes the
\texttt{\_dered} table.

\begin{warnbox}[The automatic lookup needs astroquery]
It is not in the base install or in either conda environment. Either
install it (\texttt{pip install -e ".[dered]"}), or skip the lookup by
giving the value yourself with \texttt{-{}-mw-ebv 0.0158}. Without
either you get \texttt{auto E(B-V) lookup needs astroquery (pip install
astroquery); or pass an explicit value with -{}-mw-ebv}.
\end{warnbox}

\begin{warnbox}[Use -{}-build with -{}-deredden]
\texttt{-{}-deredden} tells the fit to look for the dereddened table. If
you do not also pass \texttt{-{}-build}, that table may not exist and the
fit will tell you so.
\end{warnbox}

\subsection{Redo a fit that already exists}

By default a batch skips galaxies already fitted with the same settings.
To force it:

\begin{cmd}\ttfamily
sedfit batch -{}-roster roster.json -{}-config quick.json -{}-no-resume
\end{cmd}

If it then complains that the existing run came from different software,
add \texttt{-{}-force}.

\subsection{Use more or fewer processors}

\begin{cmd}\ttfamily
sedfit batch -{}-roster roster.json -{}-config quick.json -{}-workers 4
\end{cmd}

\texttt{-{}-workers 1} runs everything in one process, which makes error
messages much easier to read. Use it when debugging.

\subsection{Redraw figures}

\begin{cmd}\ttfamily
sedfit plot -{}-run-dir gal\_0001/SED/plain/eazy/64f456f8835f-quick\_v1
\end{cmd}

\subsection{See everything that has been run}

\begin{cmd}\ttfamily
sedfit manifest -{}-roster roster.json
\end{cmd}

\begin{cmd}\ttfamily
792 rows, 792 current runs, 0 superseded\\
\hspace*{2em}64f456f8835f~~ok~~~~~~gal\_0001/SED/plain/eazy/...~~z50=0.1043
\end{cmd}

% ==========================================================
\section{Every flag, explained}
% ==========================================================

\texttt{sedfit <verb> -{}-help} always prints the current list. This
section explains what the less obvious ones mean.

\subsection{Shared by \texttt{fit}, \texttt{batch} and \texttt{run}}

\begin{longtable}{>{\ttfamily}l p{9.4cm}}
\toprule
\textbf{\normalfont Flag} & \textbf{Meaning} \\
\midrule
\endhead
-{}-force & Replace an existing run even though it was produced by
  different package versions. Without this, the code refuses, because the
  old and new results would not be comparable. \\
-{}-allow-registry-drift & Record a warning instead of stopping when a
  filter curve has changed since the table was built. \\
-{}-deredden & Fit the Milky-Way-corrected (\texttt{\_dered}) table. On
  \texttt{batch -{}-build} and on \texttt{run}, it also makes the build
  itself dereddening, so one flag both produces and fits the corrected
  tables. \\
-{}-no-plots & Skip the figures. Faster for large sweeps. \\
\bottomrule
\end{longtable}

\subsection{\texttt{roster}}

\begin{longtable}{>{\ttfamily}l p{9.4cm}}
\toprule
\textbf{\normalfont Flag} & \textbf{Meaning} \\
\midrule
\endhead
-{}-catalog & Your sample catalog CSV. Required. \\
-{}-campaign & Your campaign config JSON. Required. \\
-{}-out & Where to write the roster. Defaults to \texttt{roster.json}
  beside the campaign file. \\
-{}-dry-run & Report what would be declared; write nothing. \\
\bottomrule
\end{longtable}

\subsection{\texttt{build}}

\begin{longtable}{>{\ttfamily}l p{9.4cm}}
\toprule
\textbf{\normalfont Flag} & \textbf{Meaning} \\
\midrule
\endhead
-{}-roster & The roster JSON. Required. \\
-{}-target & Build one galaxy; omit for all of them. \\
-{}-recipe & Build one recipe; omit for every applicable one. \\
-{}-dry-run & Validate every pair; write nothing. \\
-{}-deredden & Write the Milky-Way-corrected \texttt{\_dered} table
  \emph{instead of} the as-measured one. One build writes one table, so
  run \texttt{build} twice --- once with the flag, once without --- if
  you want both on disk. \\
-{}-mw-ebv & Use one E(B--V) for every galaxy instead of looking each one
  up. Needs \texttt{-{}-deredden}. \\
\bottomrule
\end{longtable}

\subsection{\texttt{fit}}

\begin{longtable}{>{\ttfamily}l p{9.4cm}}
\toprule
\textbf{\normalfont Flag} & \textbf{Meaning} \\
\midrule
\endhead
-{}-roster & The roster JSON. Required. \\
-{}-target & Which galaxy. \\
-{}-recipe & Which recipe's built table to fit. \\
-{}-config & The fit config JSON. \\
-{}-label & A suffix for the output directory name. Defaults to the
  \texttt{name} in the fit config. \\
-{}-jobs & A JSON list of \{target, recipe, config\} objects run in
  order, instead of naming one. Each may also carry \texttt{label} and
  \texttt{dered} (\texttt{true}/\texttt{false}, overriding
  \texttt{-{}-deredden} for that job). Cannot be combined with
  \texttt{-{}-label} or \texttt{-{}-phot-csv}, which apply to a single
  job. \\
-{}-phot-csv & Fit this photometry file instead of the built one. For
  one-off tests. \\
-{}-dry-run & Resolve and validate; write nothing. \\
\bottomrule
\end{longtable}

\noindent \texttt{fit} needs all three of \texttt{-{}-target},
\texttt{-{}-recipe} and \texttt{-{}-config} together, or
\texttt{-{}-jobs} instead. Otherwise it stops with \texttt{fit needs
-{}-target/-{}-recipe/-{}-config, or -{}-jobs}.

\subsection{\texttt{batch} and \texttt{run}}

\begin{longtable}{>{\ttfamily}l p{9.4cm}}
\toprule
\textbf{\normalfont Flag} & \textbf{Meaning} \\
\midrule
\endhead
-{}-roster & \emph{(batch only)} The roster JSON. Required.
  \texttt{run} writes its own and does not take this. \\
-{}-config & The fit config applied to every galaxy. Required. \\
-{}-targets & A CSV whose \texttt{name} column picks the galaxies. A
  \texttt{target} column works too. \\
-{}-target & A single galaxy, instead of \texttt{-{}-targets}. Passing
  both is an error. \\
-{}-recipe & One recipe; omit for every applicable one. \\
-{}-label & A suffix for every run directory. Defaults to the
  \texttt{name} in the fit config. \\
-{}-workers & How many processes. Omit to choose automatically: up to 8
  for eazy (one fewer than your CPUs), always 1 for Prospector. Never
  more than the number of jobs. \\
-{}-build & \emph{(batch only)} Rebuild each table before fitting it.
  \texttt{run} always builds and has no such flag. \\
-{}-no-resume & \emph{(batch only)} Re-run galaxies already finished.
  \texttt{run} always resumes and has no such flag. \\
-{}-mw-ebv & One E(B--V) for every galaxy instead of the per-galaxy
  lookup. On \texttt{batch} it needs both \texttt{-{}-build} and
  \texttt{-{}-deredden}; on \texttt{run}, just \texttt{-{}-deredden},
  since \texttt{run} always builds. \\
-{}-report & Where to write the report CSV. Defaults to
  \texttt{batch\_report.csv} beside the roster. \\
-{}-log-dir & Where to write the per-job logs
  (\texttt{<target>\_\_<recipe>.log}). Defaults to
  \texttt{batch\_logs/} beside the report. \\
-{}-stop-after-failures & Give up after this many failures.
  \texttt{0} disables the limit. Default 20. \\
\bottomrule
\end{longtable}

\texttt{run} additionally takes \texttt{-{}-catalog}, \texttt{-{}-campaign}
and \texttt{-{}-out}, which behave exactly as they do for
\texttt{roster}. It does \textbf{not} take \texttt{-{}-roster},
\texttt{-{}-build}, \texttt{-{}-no-resume} or \texttt{-{}-dry-run}.

A galaxy named in \texttt{-{}-targets} that the roster does not declare
is reported as \texttt{skipped} in the report rather than failing the
batch --- unless none of them is in the roster, which stops before any
job runs.

% ==========================================================
\section{What comes out}
% ==========================================================

\subsection{The directory tree}

\begin{cmd}\ttfamily
gal\_0001/\\
\hspace*{2em}Photometry/\\
\hspace*{4em}gal\_0001\_catalog.csv\hspace*{2em}(yours, untouched)\\
\hspace*{4em}gal\_0001\_sed\_plain.csv\hspace*{1em}the combined table\\
\hspace*{4em}gal\_0001\_sed\_plain.provenance.json\\
\hspace*{2em}SED/\\
\hspace*{4em}plain/\hspace*{9em}one folder per recipe\\
\hspace*{6em}eazy/\hspace*{8.5em}one folder per backend\\
\hspace*{8em}64f456f8835f-quick\_v1/\hspace*{1em}one run\\
\hspace*{10em}config.json\\
\hspace*{10em}phot.csv\\
\hspace*{10em}phot.provenance.json\\
\hspace*{10em}manifest.json\\
\hspace*{10em}summary.csv\\
\hspace*{10em}arrays.npz\\
\hspace*{10em}plots/
\end{cmd}

\noindent An eazy run writes a few more files than these --- the filter
and template records it used. A Prospector run writes
\texttt{result.h5} instead of \texttt{summary.csv} and
\texttt{arrays.npz}.

The long name on the run folder is not random. It is a fingerprint of the
settings, the photometry and the filter curves. Fit the same galaxy the
same way twice and you get the same folder --- it is replaced, not
duplicated. Change anything that matters and you get a new folder, so
results from different settings can never be confused.

\subsection{What each file is}

\begin{center}
\begin{tabular}{>{\ttfamily}l p{9.2cm}}
\toprule
\textbf{\normalfont File} & \textbf{What it holds} \\
\midrule
\_sed\_<recipe>.csv & The combined photometry the fit used \\
.provenance.json & Every input file, hashed, and every choice made \\
config.json & The complete settings, with defaults filled in \\
phot.csv & An exact copy of the photometry this run fitted \\
manifest.json & This run's result line \\
summary.csv & Redshifts and chi-square (eazy only) \\
arrays.npz & The redshift grid and the probability curve (eazy only) \\
result.h5 & The chain, model stamps and run parameters (Prospector only) \\
plots/ & The figures \\
batch\_report.csv & One row per job --- one galaxy under one recipe \\
runs.jsonl & Every run ever, appended \\
\bottomrule
\end{tabular}
\end{center}

\subsection{The numbers you probably want}

\texttt{batch\_report.csv} is the quickest place to look:

\begin{center}
\begin{tabular}{>{\ttfamily}l p{8.8cm}}
\toprule
\textbf{\normalfont Column} & \textbf{Meaning} \\
\midrule
target, recipe & which galaxy, which recipe \\
status & \texttt{ok}, \texttt{failed}, or \texttt{skipped} \\
stage & which phase a job died in: \texttt{roster},
  \texttt{applicable}, \texttt{build}, \texttt{plan}, \texttt{fit} \\
zred\_p50 & the redshift estimate (the median) \\
zred\_p16, zred\_p84 & the 1-sigma range \\
n\_bands\_fit & how many bands actually entered the fit \\
n\_bands\_table & how many were available \\
run\_id, path & where the run directory is \\
seconds & how long it took \\
error & why it failed, when it did \\
\bottomrule
\end{tabular}
\end{center}

If \texttt{n\_bands\_fit} is much smaller than \texttt{n\_bands\_table},
bands were dropped --- usually for low signal-to-noise. The run's
\texttt{manifest.json} lists exactly which ones and why. When a job
fails, read \texttt{stage} first: it says which phase gave out, which
usually tells you whether the problem is in your files or in the fit.

\begin{notebox}[The same numbers have two names]
\texttt{summary.csv} calls the percentiles \texttt{z160},
\texttt{z500} and \texttt{z840}. The batch report and
\texttt{manifest.json} call the identical numbers \texttt{zred\_p16},
\texttt{zred\_p50} and \texttt{zred\_p84}. Searching a
\texttt{summary.csv} for \texttt{zred\_p50} finds nothing --- look for
\texttt{z500}.
\end{notebox}

% ==========================================================
\section{Reading the figures}
% ==========================================================

An eazy run writes two figures into \texttt{plots/} --- three when
\texttt{z\_fixed} is set, which adds the SED at that fixed redshift. A
Prospector run writes three different ones: a corner plot of the
posterior, the sampler trace, and the SED at the best-fit parameters.
The two described below are the eazy pair.

\textbf{The SED figure} shows your photometry as points, with error bars,
against the best-fitting template. Points are colored by instrument. If
the template misses a group of points badly, that instrument may be on
the wrong flux scale --- check the recipe.

\textbf{The redshift-scan figure} shows how well each redshift fits. A
good fit shows one clear, narrow minimum. Two things to watch for:

\begin{itemize}
  \item \textbf{Several comparable minima} --- the fit cannot decide, and
        the quoted uncertainty understates the real one.
  \item \textbf{The minimum at the edge of the range} --- the true
        redshift is probably outside your \texttt{z\_min}--\texttt{z\_max}
        window. The code warns about this in the log too. Widen the range
        and re-run.
\end{itemize}

% ==========================================================
\section{When something goes wrong}
% ==========================================================

Error messages name the file and say what to do. The common ones:

\vspace{4pt}
\noindent\textbf{\texttt{roster not found: ...}}\\
The path after \texttt{-{}-roster} is wrong, or you have not run
\texttt{sedfit roster} yet.

\vspace{6pt}
\noindent\textbf{\texttt{unknown target 'x'; the roster declares ...}}\\
A typo, or the galaxy was dropped during roster generation. Check
\texttt{roster.json.report.csv} for why.

\vspace{6pt}
\noindent\textbf{\texttt{... has not been built; run: sedfit build ...}}\\
Exactly what it says --- the message includes the command to run.

\vspace{6pt}
\noindent\textbf{\texttt{z\_ref\_kind 'reference' requires the roster to
declare reference\_redshift}}\\
A catalog row says \texttt{reference} but the campaign config has no
\texttt{reference\_redshift} for it to adopt. Add it to the campaign and
regenerate the roster. Note this fires on the verb \emph{after}
\texttt{roster}, because generation does not check it.

\vspace{6pt}
\noindent\textbf{\texttt{sedfit: command not found}}\\
The environment exists but the package is not installed into it. From
the repository folder: \texttt{conda activate <env>} then \texttt{pip
install -e .}. The \texttt{.yml} files do not do this for you.

\vspace{6pt}
\noindent\textbf{\texttt{No module named 'pytest'}}\\
The base install does not include it. \texttt{pip install -e ".[test]"}.

\vspace{6pt}
\noindent\textbf{\texttt{auto E(B-V) lookup needs astroquery ...}}\\
\texttt{-{}-deredden} without \texttt{-{}-mw-ebv} needs
\texttt{astroquery}. Either \texttt{pip install -e ".[dered]"} or supply
the value yourself with \texttt{-{}-mw-ebv}.

\vspace{6pt}
\noindent\textbf{\texttt{unknown bands: [...]}}\\
A band name in the campaign config is not one the package knows. The
legal names are in \texttt{sedfit/data/registry.json}.

\vspace{6pt}
\noindent\textbf{\texttt{unknown provider token 'x' (known: [...])}}\\
A \texttt{provider} in the campaign config is not one of the ten legal
tokens. The message lists them.

\vspace{6pt}
\noindent\textbf{\texttt{unknown recipe 'x'; the roster declares [...]}}
or \textbf{\texttt{recipes not in the roster: [...]}}\\
The name after \texttt{-{}-recipe} is not one the campaign defines. The
message lists what is available.

\vspace{6pt}
\noindent\textbf{\texttt{templates=... is not a directory, a .param file,
or a packaged set}}\\
The message lists the available set names. Check the spelling in your fit
config.

\vspace{6pt}
\noindent\textbf{\texttt{only N fittable bands after inclusion and gates
(min\_valid\_bands=5)}}\\
Too few bands survived. Either the photometry is thin for that galaxy, or
the signal-to-noise cut removed too much. Lower
\texttt{min\_valid\_bands} if that is scientifically appropriate --- it
sits at the \textbf{top level} of the fit config, beside
\texttt{err\_floor}, not inside the \texttt{eazy} block. Its default is
5; \texttt{min\_snr\_broadband} defaults to 2.0 and \texttt{err\_floor}
to 0.05.

\vspace{6pt}
\noindent\textbf{\texttt{WARNING: 'gal\_0001': chi2 minimum on the grid
edge (z=0.6000); widen the grid}}\\
Increase \texttt{z\_max} or decrease \texttt{z\_min}. This is a warning,
not an error: the run completes and writes results, but the redshift is
pinned against the edge of the window you gave it and should not be
trusted. Worth checking whenever you have not set the grid yourself ---
the defaults are 0.05 to 0.16.

\vspace{6pt}
\noindent\textbf{\texttt{existing run was produced by different machinery}}\\
You reinstalled or upgraded something since that run. Add
\texttt{-{}-force} to replace it.

\vspace{6pt}
\noindent\textbf{\texttt{no declared band}} in the roster report\\
The file exists but holds none of the bands the campaign config expects
from it. Usually the \texttt{provider} or the band names are wrong.

\vspace{6pt}
\noindent\textbf{A command seems to hang.}\\
The first step of every verb opens every photometry file the roster
declares. On a network drive or a cloud-synced folder, that can take
minutes the first time. It is faster afterwards.

\vspace{6pt}
\noindent\textbf{An unreadable error during a batch.}\\
Re-run the failing galaxy on its own with \texttt{-{}-workers 1}:

\begin{cmd}\ttfamily
sedfit fit -{}-roster roster.json -{}-target gal\_0007 \textbackslash\\
\hspace*{4em}-{}-recipe plain -{}-config quick.json
\end{cmd}

Per-job logs from a batch are in \texttt{batch\_logs/}, one file per
galaxy-and-recipe, named \texttt{<target>\_\_<recipe>.log}.

% ==========================================================
\section{Getting more out of it}
% ==========================================================

\subsection{Several recipes at once}

A recipe is one way of combining your photometry. Define several in the
campaign config and \texttt{batch} runs them all, giving each its own
folder. Comparing recipes is how you test whether a result depends on how
the instruments were combined.

A recipe's \texttt{reference} says what defines the flux scale, and each
source takes a \texttt{role}:

\begin{center}
\begin{tabular}{>{\ttfamily}l p{8.6cm}}
\toprule
\textbf{\normalfont reference} & \textbf{What is truth} \\
\midrule
spherex & the SPHEREx spectrum; overlapping broadbands are rescaled onto it \\
anchors & one or two named broadband groups; the spectrum is moved onto them \\
none & nothing; every source stays at its measured level \\
\bottomrule
\end{tabular}
\end{center}

\begin{center}
\begin{tabular}{>{\ttfamily}l p{8.6cm}}
\toprule
\textbf{\normalfont role} & \textbf{What the source does} \\
\midrule
anchor & defines the reference. Legal only under \texttt{anchors}, which
  needs one or two. One rescales the spectrum without changing its
  color; two tilt it. \\
stitch & its whole instrument group is multiplied by one factor,
  measured in the reddest band it shares with the reference. Legal under
  \texttt{spherex} and \texttt{anchors}. \\
float & carried through untouched. Legal everywhere, and the only role
  legal under \texttt{none}. \\
\bottomrule
\end{tabular}
\end{center}

A source can only stitch if it actually overlaps the reference; the
build says so rather than guessing. \texttt{examples/README.md} works
through four recipes covering every combination.

\subsection{Comparing template sets}

Copy your fit config, change \texttt{templates} and \texttt{name}, and run
both. Because the run folder is named after the settings, the two results
sit side by side without overwriting each other. The shipped sets are:

\begin{center}
\begin{tabular}{>{\ttfamily}l p{8.2cm}}
\toprule
\textbf{\normalfont Set} & \textbf{What it is} \\
\midrule
brown14\_vac\_cosmos160 & The recommended set: Brown+14 galaxies plus the
  COSMOS templates \\
brown14 & The Brown+14 atlas as published \\
brown14\_vac & Brown+14, air-to-vacuum corrected \\
ssp\_miles & FSPS simple stellar populations, MILES library \\
ssp\_c3k\_a & The same on the C3K library \\
\bottomrule
\end{tabular}
\end{center}

\subsection{The central log}

Every finished run appends one line to \texttt{runs.jsonl} under your data
root. It is one JSON object per line, so it loads directly into pandas:

\begin{cmd}\ttfamily
import pandas as pd\\
runs = pd.read\_json("sed\_fitting/runs.jsonl", lines=True)
\end{cmd}

That is the easiest way to compare many runs at once.

\subsection{A worked example to copy}

The repository ships a complete, runnable example in \texttt{examples/}:
three synthetic galaxies with broadband photometry and SPHEREx spectra,
the input files with every field annotated (including a Prospector
config), and four recipes covering every way the package can combine
instruments. Run it as-is:

\begin{cmd}\ttfamily
cd examples\\
sedfit run -{}-catalog sample\_catalog.csv -{}-campaign campaign.json \textbackslash\\
\hspace*{4em}-{}-config fit\_quick.json
\end{cmd}

Then copy \texttt{sample\_catalog.csv}, \texttt{campaign.json} and
\texttt{fit\_quick.json} into your own working directory and replace
their contents. \texttt{examples/README.md} explains every field.

\subsection{Where to read next}

\texttt{docs/technical\_manual.md} in the repository explains how the code
works internally --- what each module does, how the flux scaling is
computed, and what every output field means.

\end{document}
